% Clase 15 --- El ciclotrón (§2.2).
% El docente hizo dos dibujos "porque me cuesta hacer todo en un solo dibujo":
% la vista de las dos Ds con la ranura, y el detalle del movimiento. Acá van los
% dos, con el énfasis en dónde se gana energía (sólo en la ranura) y en por qué
% alcanza UNA sola frecuencia para la fuente alterna.
\begin{center}
\begin{tikzpicture}[scale=1]

  % =============== (a) las dos Ds y la espiral ===============
  \begin{scope}
    % D izquierda
    \draw[figgray, line width=1.1pt, fill=figgray!10]
      (-0.12,2.0) arc (90:270:2.0) -- cycle;
    % D derecha
    \draw[figgray, line width=1.1pt, fill=figgray!10]
      (0.12,-2.0) arc (-90:90:2.0) -- cycle;
    \node[figlbl] at (-1.15,1.35) {D};
    \node[figlbl] at (1.15,1.35) {D};

    % campo perpendicular
    \foreach \x in {-1.55,-0.8,0.8,1.55} {
      \foreach \y in {-1.45,-0.7,0.05,0.8} {
        \node[figblue!45, scale=0.55] at (\x,\y) {$\otimes$};
      }
    }

    % la espiral: semicírculos de radio creciente
    \draw[figred, line width=1pt] (0,0.3) arc (90:-90:0.3);
    \draw[figred, line width=1pt] (0,-0.3) arc (-90:-270:0.6);
    \draw[figred, line width=1pt] (0,0.9) arc (90:-90:0.9);
    \draw[figred, line width=1pt] (0,-0.9) arc (-90:-270:1.2);
    \draw[figred, line width=1pt] (0,1.5) arc (90:-90:1.5);
    \node[figcarga, fill=figred, minimum size=5pt] at (0,0.3) {};

    % salida
    \draw[figvecres] (0.1,-1.5) -- (1.0,-2.35);
    \node[figlbl, figred, anchor=west] at (1.05,-2.4) {sale};

    \node[figlbl, anchor=north, align=center] at (0,-3.0)
      {(a) vista desde arriba};
  \end{scope}

  % =============== (b) qué pasa en la ranura ===============
  \begin{scope}[xshift=7.0cm]
    \fill[figgray!10] (-2.0,-1.4) rectangle (-0.25,1.4);
    \fill[figgray!10] (0.25,-1.4) rectangle (2.0,1.4);
    \draw[figgray, line width=1.1pt] (-2.0,-1.4) rectangle (-0.25,1.4);
    \draw[figgray, line width=1.1pt] (0.25,-1.4) rectangle (2.0,1.4);
    \node[figlblaux, anchor=south] at (0,1.5) {la ranura};

    % la diferencia de potencial alterna
    \draw[figvecres, line width=1.2pt] (-0.2,0.45) -- (0.2,0.45);
    \draw[figvecaux, figblue, line width=1.1pt,
          -{Latex[length=2.2mm,width=1.6mm]}] (0.2,-0.45) -- (-0.2,-0.45);

    \node[figlbl, anchor=center, align=center] at (-1.12,0)
      {adentro:\\[-0.2em] \footnotesize sólo gira};
    \node[figlbl, anchor=center, align=center] at (1.12,0)
      {adentro:\\[-0.2em] \footnotesize sólo gira};

    \node[figlbl, figred, anchor=north, align=center] at (0,-1.6)
      {la energía se gana \emph{sólo} acá:\\[-0.2em]
       \footnotesize un $\Delta V$ que cambia de signo cada media vuelta};

    \node[figlbl, anchor=north, align=center] at (0,-3.0)
      {(b) detalle de la ranura};
  \end{scope}

\end{tikzpicture}\\[0.5em]
{\small\itshape Adentro de cada D no hay campo eléctrico: la partícula sólo
describe su semicírculo, y como la fuerza magnética no hace trabajo, entra y
sale con la misma rapidez. Toda la energía se entrega en la ranura, y para que
la partícula sea acelerada ---y no frenada--- al llegar allí, la diferencia de
potencial tiene que haber cambiado de signo en el ínterin: de ahí que la fuente
sea alterna. Acá es donde se cobra el resultado de §1.3: como el período no
depende de la velocidad, \emph{todas} las medias vueltas duran lo mismo aunque
la partícula vaya cada vez más rápido, así que basta una \textbf{única
frecuencia fija} para acompañarla durante todo el recorrido. Los semicírculos
van creciendo hasta que la partícula alcanza el borde y sale con
$E_c = q^2B^2R^2/2m$. Las dos limitaciones salen de ahí: la energía crece con
$R^2$ ---y el imán no puede ser arbitrariamente grande--- y, a velocidades
relativistas, el período deja de ser constante y el diseño entero se cae. Por
eso los aceleradores de kilómetros usan otra tecnología, y los ciclotrones
quedaron para medicina y estudio de materiales.}
\end{center}
